%قبل از شروع مسئله ترجیحاً چیزی ننویسید. اما اگر خواستید بنویسید، در قالب لیستی (فقرات و شمارش) بنویسید.
%برای اضافه کردن نمره، بعدی از شروع‌مسئله، [x] را اضافه کنید و به جای x نمره مورد نظر را قرار دهید. برای زیرمسئله‌ها هم به همین صورت است.
%\شروع‌مسئله[10]
\شروع‌مسئله [15]

یکی از قدرت‌های ماوراء الطبیعه‌ی شنگدباو قدرت میانه‌روی اوست! او به قدری میانه‌رو است که آقاصادق یک روز به او کلی عدد به ترتیب می‌گوید و پس از هر بار شنیدن یک عدد جدید شنگدباو میانه‌‌ی همه‌ی اعدادی که تا آن‌جا شنیده را به سرعت اعلام می‌کند!

میانه‌ی $n$ عدد برابر عدد $\lceil \frac{n}{2} \rceil$ام در ترتیب مرتب شده‌شان است. یعنی برای مثال میانه‌ی سه عدد ۱، ۱۰ و ۵ برابر ۵ می‌باشد. یا میانه‌ی چهار عدد ۱، ۱۰، ۱۱، ۱۰۰ برابر ۱۰ می‌باشد. 

آیا می‌توانید مثل شنگدباو باشید؟

%\متن‌سیاه{سوال عملی اول:}\\
%\متن‌سیاه{دراپ‌تست}\\

\متن‌سیاه{ورودی:}\\
در ورودی ابتدا به‌شما عدد $q$ داده‌ می‌شود که تعداد اعدادی است که آقاصادق به شنگدباو می‌گوید و سپس در $q$ خط بعدی در هر خط یک عدد $x$ می‌آید که عددی است که آقاصادق در آن مرحله اعلام کرده‌است.

%در خط اول ورودی، ابتدا k و سپس m داده می‌شود.\\

\متن‌سیاه{خروجی:}\\
در $q$ خط خروجی دهید و در خط $i$ام میانه‌ی اعداد اول تا $i$ام را اعلام کنید.
%در تنها سطر خروجی، تعداد آیفون‌های مورد نیاز می‌آید. اگر در هیچ صورتی چنین کاری ممکن نبود ۱- چاپ می‌شود.\\

\متن‌سیاه{محدودیت‌ها:}
$$ 1 \le q \le 10^5 \ , \ 1 \le x \le 10^9$$

%$m\leq10^9$
%\\
\متن‌سیاه{ورودی و خروجی نمونه}\\
%% تست اول
\begin{latin}
 \begin{tabularx}{\textwidth}{|X | X|}
 \hline
 \textbf{stdin} &
 \textbf{stdout}\\
 \hline
 5
 1
 10
 3
 4
 9&
 1
 1
 3
 3
 4\\
 \hline
 
 \end{tabularx}
\end{latin}

%% تست دوم
%
%\begin{latin}
% \begin{tabularx}{\textwidth}{|X | X|}
% \hline
% \textbf{stdin} &
% \textbf{stdout}\\
% \hline
% 340 20&
% 3\\
% \hline
% 
% \end{tabularx}
%\end{latin}}
% صورت اصلی مسئله را در این‌جا بنویسید.
%اگر مسئله چند قسمت داشت به این صورت عمل کنید:
%\شروع{زیرمسائل}
%\زیرمسئله[0] % مسئلۀ ۱آ
%\زیرمسئله % مسئلۀ ۱ب
%\زیرمسئله % مسئلۀ ۱پ
%\زیرمسئله % مسئلۀ ۱ت
%\پایان{زیرمسائل}

\solution{
دو تا درخت هرمی استفاده می‌کنیم یکی درخت هرمی کمینه از عناصر بعد از میانه‌ و دیگری درخت هرمی بیشینه از عناصر قبل از میانه‌است (شامل خود میانه‌). ساختن درخت هرمی بیشینه همانند همان درخت هرمی کمینه‌است فقط هر رأس از بچه‌هایش بزرگتر است و ریشه بیشینه می‌باشد. 

در این‌صورت همیشه بزرگترین عنصر هرم بیشینه برابر جواب است. هرگاه عددی اضافه می‌کنیم، آن‌را به هرم بیشینه اضافه می‌کنیم و اگر لازم بود بزرگترین عنصر هرم بیشینه را حذف کرده و به هرم کمینه‌ اضافه می‌کنیم. به این‌صورت همیشه در این دو هرم اعداد نیمه‌ی اول و دوم به صورت مجزا باقی می‌مانند. 

راه‌حل ‌$O(qlogq)$ است.  

}




%%%%%%%%%%%%%%%%%%%%%%%%%%%%%%%%%%%%%%%%%%%%%%%%%%%%%
% برخی دستورات مفید برای تمرین‌ها. %
%%%%%%%%%%%%%%%%%%%%%%%%%%%%%%%%%%%%%%%%%%%%%%%%%%%%%

% فرمول‌های نمونه:
%$f(x) = \Theta(x)$
%$T(x, y) \leq \log(x) + 2^y + \binom{2n}{n}$
% {\tt code\_function}


% لیست موارد بدون شماره‌گذاری:
%\شروع{فقرات}
% \فقره مورد اول 
% \شروع{فقرات}
% \فقره مورد اول داخلی 
% \شروع{فقرات}
% \فقره اولین مورد خیلی داخلی 
% \فقره دومین مورد خیلی داخلی 
% \پایان{فقرات}
% \فقره مورد دوم داخلی 
% \پایان{فقرات}
% \فقره مورد دوم 
%\پایان{فقرات}

% لیست موارد شماره‌گذاری شده:
%\شروع{شمارش}
% \فقره یک مورد 
% \فقره یک مورد دیگر 
% \فقره باز هم یک مورد دیگر
%\پایان{شمارش}

% برای نوشتن یک فرمول چند خطی که یک شماره بخوره. برای حالت‌های دیگر به سایت زیر مراجعه کنید:
% https://www.overleaf.com/learn/latex/Aligning_equations_with_amsmath
%\begin{align} 
% \begin{split}
% (x+y)^3 &= (x+y)^2(x+y) \\
% &= (x^2+2xy+y^2)(x+y) \\
% &= (x^3+2x^2y+xy^2) + (x^2y+2xy^2+y^3) \\
% &= x^3+3x^2y+3xy^2+y^3
% \end{split} 
%\end{align}

% نوشتن ماتریس و آرایه دوبعدی
%\begin{align}
% A = 
% \begin{bmatrix}
% A_{11} & A_{21} \\
% A_{21} & A_{22}
% \end{bmatrix}
%\end{align}

% اضافه کردن تصویر با شرح. برای اضافه کردن تصویر بدون شرح، فقط دستور درج‌تصویر را استفاده کنید.
%\شروع{شکل}[h!]
%\تنظیم‌ازوسط
% \درج‌تصویر[پهنا=0.5\پهنای‌متن]{img.jpg}
%\شرح{شرح تصویر}
%\برچسب{تصویر:۱}
%\پایان{شکل}
%\پایان{زیرمسائل}

%کد زیر را ببینید.
%%\lstinputlisting{source_filename.py}
%\begin{latin}
%\begin{lstlisting}
%# Import the modules
%import sys
%import random
%
%ans = True
%
%while ans:
%question = raw_input("Ask the magic 8 ball a question: (press enter to quit) ")
%
%answers = random.randint(1,8)
%
%if question == "":
%sys.exit()
%
%elif answers == 1:
%print "It is certain"
%
%elif answers == 2:
%print "Outlook good"
%
%elif answers == 3:
%print "You may rely on it"
%
%elif answers == 4:
%print "Ask again later"
%
%elif answers == 5:
%print "Concentrate and ask again"
%
%elif answers == 6:
%print "Reply hazy, try again"
%
%elif answers == 7:
%print "My reply is no"
%
%elif answers == 8:
%print "My sources say no"
%\end{lstlisting}
%\end{latin}